% =============================================================================
% STAGING EXTRACTION FOR SPACE SEGMENT DESIGN DEFINITION FILE (DDF)
% Target Document : EWAS-HSB-SSEG-DDF-03 (ECSS-E-ST-10C Annex G)
% Source Document : EWAS-HSB-SSEG-TRS-03 (sections/04_selected_concept.tex)
% Extracted Scope : Section 4.2 (OISL Terminal Hardware & Routing),
%                   Section 4.3 (Optical Payload & LOS Steering Design),
%                   Section 4.4 (Product Breakdown, WBS, Careful COTS Bus, Table 6)
% Status          : Staging file only; verbatim design definitions preserved
%                   prior to streamlining the Level-2 Space Segment TRS.
% =============================================================================
\documentclass[11pt,a4paper]{article}

\usepackage[utf8]{inputenc}
\usepackage[T1]{fontenc}
\usepackage{lmodern}
\usepackage[a4paper, top=2.5cm, bottom=2.5cm, left=2.5cm, right=2.5cm]{geometry}
\usepackage{longtable}
\usepackage[table]{xcolor}
\usepackage{tabularx}
\usepackage{array}
\usepackage{float}
\usepackage{amsmath}
\usepackage[hidelinks]{hyperref}

\definecolor{TitleBlue}{HTML}{1F497D}
\definecolor{HeadingBlue}{HTML}{365F91}
\definecolor{SubHeadingBlue}{HTML}{4F81BD}
\definecolor{TableHeaderBlue}{HTML}{4F81BD}
\definecolor{TableRowLightBlue}{HTML}{DBE5F1}
\definecolor{TableBorderBlue}{HTML}{95B3D7}

\setlength{\parindent}{0pt}
\setlength{\parskip}{0.6em}
\setcounter{secnumdepth}{-1}

\makeatletter
\renewcommand\section{\@startsection{section}{1}{\z@}%
    {-3.0ex \@plus -1ex \@minus -.2ex}%
    {1.5ex \@plus .2ex}%
    {\normalfont\Large\bfseries\sffamily\color{HeadingBlue}}}

\renewcommand\subsection{\@startsection{subsection}{2}{\z@}%
    {-2.5ex \@plus -1ex \@minus -.2ex}%
    {1.0ex \@plus .2ex}%
    {\normalfont\large\bfseries\sffamily\color{SubHeadingBlue}}}

\renewcommand\subsubsection{\@startsection{subsubsection}{3}{\z@}%
    {-2.0ex \@plus -1ex \@minus -.2ex}%
    {0.8ex \@plus .2ex}%
    {\normalfont\normalsize\bfseries\sffamily\color{SubHeadingBlue}}}
\makeatother

\newcommand{\thead}[1]{\textcolor{white}{\textbf{\textsf{#1}}}}
\providecommand{\abbr}[1]{#1}

\begin{document}

\section{EWAS-HSB-SSEG-DDF-03 Staging Extraction (from EWAS-HSB-SSEG-TRS-03)}

This standalone staging file preserves all detailed hardware, optical, mechanical, communications, and product-tree design definitions extracted verbatim from Section~4 of \texttt{EWAS-HSB-SSEG-TRS-03} (\texttt{sections/04\_selected\_concept.tex}) for incorporation into the Space Segment Design Definition File (\texttt{EWAS-HSB-SSEG-DDF-03}, ECSS-E-ST-10C Annex~G).

% =============================================================================
% EXTRACTED FROM TRS SECTION 4.2: OISL TERMINAL HARDWARE AND ROUTING
% =============================================================================
\subsubsection{4.2 Mandatory Space Segment Optical Inter-Satellite Link (OISL) Network and Routing Concept}

A core architectural pillar of the \abbr{EWAS} Space Segment is its \textbf{mandatory laser Optical Inter-Satellite Link (\abbr{OISL}) backbone} (\texttt{R-023}, \texttt{R-031}--\texttt{R-033}, \texttt{R-044}). In an extra-regional launch scenario (\texttt{R-049}), the two \abbr{LEO} satellites best positioned to observe a boost or post-boost missile plume may be located over continental Eurasia, the Arctic, or oceanic regions outside the direct radio-frequency line-of-sight of the European Ground Station Network (\abbr{WBS} \texttt{123456.02.05}). Whereas Space-to-Ground links (\texttt{IF-EXT-01}) use Radio Frequency (\abbr{RF}: S/X/Ka-band) for all-weather cloud penetration, all Space-to-Space links operate above the atmosphere and use 1,550~nm laser \abbr{OISL} terminals (backed up only by an emergency low-rate S-band omni beacon in Safe Mode):
\begin{enumerate}
    \item \textbf{Intra-Layer \abbr{LEO} Optical Inter-Satellite Link (\abbr{OISL}) Mesh:} Each of the 72 \texttt{SATL} spacecraft carries \textbf{four (4) laser \abbr{OISL} terminals} operating in the 1,550~nm (C-band) optical window. Each terminal delivers a full-duplex user data rate of \textbf{$\ge 500.0\text{ Mbps}$} (sized to transport the uncompressed stream of two $512 \times 512$ \abbr{SWIR}/\abbr{MWIR} sensors at 30~Hz multi-bit depth plus $25\%$ framing and relay overhead) over inter-satellite link distances up to \textbf{$\ge 6{,}500\text{ km}$}:
    \begin{itemize}
        \item \emph{Two Permanent Intra-Plane Terminals (Fore and Aft):} Maintain continuous optical lock with the immediately preceding and trailing satellites within the same orbital plane ($\Delta u = 40^\circ$, constant chord distance of 5,047~km, zero relative azimuth rate). This forms 8 resilient longitudinal rings around the Earth.
        \item \emph{Two Agile Cross-Plane Terminals (Left and Right):} Use coarse-pointing hemispherical periscopes and fine-steering mirrors to maintain optical links across adjacent orbital planes ($\Delta\Omega = 45^\circ$). At the equator ($a = 7{,}378.14\text{ km}$ with $\Delta u_{\text{inter}} = 22.5^\circ$), the cross-plane chord distance reaches 5,647--6,144~km while keeping the optical line of sight $> 245\text{ km}$ above Earth's surface (well clear of atmospheric refraction and cloud layers). Specifying a maximum \abbr{OISL} operating range of $\ge 6{,}500\text{ km}$ guarantees continuous cross-plane connectivity from the equator to high latitudes, dynamically re-handshaking only across the narrow $\pm 80^\circ$ polar cross-seam.
    \end{itemize}
    \item \textbf{Cross-Layer \abbr{GEO}--\abbr{LEO} Optical Relay Trunks:} To avoid multi-hop routing delays across multiple \abbr{LEO} planes, each \texttt{SATG} geostationary satellite carries agile 1,550~nm laser \abbr{OISL} relay terminals providing an aggregate data rate of \textbf{$\ge 1.5\text{--}2.0\text{ Gbps}$} (accommodating $3\text{--}4$ simultaneous \texttt{SATL} sensor streams plus command and relay overhead). Because the maximum geometric line-of-sight distance between a \abbr{GEO} satellite ($a = 42{,}164\text{ km}$) and a 1,000~km \abbr{LEO} satellite near the Earth limb reaches $\sim 45{,}195\text{ km}$, the cross-layer \abbr{OISL} terminals are qualified across an operating link range of \textbf{$34{,}800\text{ km}$ to $45{,}000\text{ km}$}, bridging the \abbr{LEO} tracking layer directly to European ground gateways.
    \item \textbf{Autonomous Software-Defined Mesh Routing and Stereoscopic Vector Exchange:} All \abbr{OISL} packets are encapsulated in \abbr{CCSDS}-compliant frames protected by hardware bulk encryption (\texttt{R-031}) and narrow-beam Low Probability of Intercept and Anti-Jamming (\abbr{LPI}/\abbr{AJ}) safeguards (\texttt{R-032}). During an active warning sequence (Table~\ref{tab:warning_chain_latency}), the \abbr{OISL} network executes three time-critical functions:
    \begin{itemize}
        \item \emph{Sub-3-Second Cueing Delivery (\texttt{R-044}):} Routes the Ground Segment's \emph{LEO Tracking Request} from the nearest ground gateway or via \texttt{SATG} cross-layer optical relay to the two selected \texttt{SATL} spacecraft within $\le 3.0\text{ s}$.
        \item \emph{Inter-Satellite Stereoscopic \abbr{LOS} Exchange (\texttt{R-023}):} Once both cued \texttt{SATL} satellites lock onto the target ($< 5\text{ s}$), they exchange GNSS-synchronized ($< 50~\mu\mathrm{s}$ epoch precision) timestamped \abbr{LOS} unit vectors and spacecraft position/attitude covariance matrices directly across the \abbr{OISL} mesh at $10\text{--}20\text{ Hz}$. This enables the primary \texttt{SATL} node's On-Board Data Handling (\abbr{OBDH}) processor to compute a real-time \emph{Onboard Preliminary Trajectory} estimate via dual-satellite stereoscopic triangulation (\texttt{R-023}).
        \item \emph{Immediate Dual-Product Downlink (\texttt{R-012}, \texttt{R-023}, \texttt{R-033}):} Simultaneously streams both the \emph{LEO Observation Packages} (from both \abbr{LEO} satellites) and the \emph{Onboard Preliminary Trajectory} estimate back to the Ground Segment Data Processing System (\texttt{123456.02.04}) within $\le 5\text{--}7\text{ s}$ for authoritative multi-sensor ground fusion (\texttt{R-026}, \texttt{R-027}, \texttt{R-045}).
    \end{itemize}
\end{enumerate}

% =============================================================================
% EXTRACTED FROM TRS SECTION 4.3: MISSION PAYLOAD & OPTICAL LOS STEERING
% =============================================================================
\subsubsection{4.3 Mission Payload Architecture and Independent Optical Payload LOS Steering Concept}

Both the \texttt{SATG} detection payload and the \texttt{SATL} tracking payload are built around cooled dual-band infrared telescopes engineered to exploit the radiometric signature of rocket exhaust plumes (\texttt{R-009}, \texttt{R-010}, \texttt{R-020}--\texttt{R-025}):
\begin{itemize}
    \item \textbf{Dual-Band \abbr{SWIR}/\abbr{MWIR} Spectral Architecture (\texttt{R-009}, \texttt{R-020}):} Solid- and liquid-propellant rocket plumes emit intense thermal and vibrational-rotational radiation dominated by hot water vapour ($\mathrm{H_2O}$) and carbon dioxide ($\mathrm{CO_2}$). Both payloads employ dichroic beam-splitters feeding cooled 2D Mercury Cadmium Telluride ($\mathrm{HgCdTe}$ / \abbr{MCT}) Focal Plane Arrays (\abbr{FPA}s) operating simultaneously in two spectral bands:
    \begin{enumerate}
        \item \emph{Short-Wave Infrared (\abbr{SWIR}, $2.65\text{--}2.85~\mu\mathrm{m}$):} Centered on the strong $\mathrm{H_2O}$ and $\mathrm{CO_2}$ emission band, providing high contrast against solar-reflected cloud backgrounds during lower- and mid-atmospheric boost combustion.
        \item \emph{Medium-Wave Infrared (\abbr{MWIR}, $4.15\text{--}4.45~\mu\mathrm{m}$):} Centered on the $\mathrm{CO_2}$ ``red/blue spike'' wing bands, enabling high-sensitivity detection and tracking as the missile ascends into the upper stratosphere and exosphere against cold space and low-temperature cloud tops.
    \end{enumerate}
    Simultaneous dual-band pixel co-registration allows the onboard processor to perform real-time two-colour spectral ratioing ($I_{\text{SWIR}}/I_{\text{MWIR}}$). This discriminates genuine rocket exhaust plumes ($T_{\text{eff}} \approx 1{,}800\text{--}3{,}000\text{ K}$) from solar specular glint ($T_{\text{eff}} \approx 5{,}800\text{ K}$) and ambient terrestrial fires ($T_{\text{eff}} \approx 800\text{--}1{,}100\text{ K}$, \texttt{R-019}, \texttt{R-022}).
    \item \textbf{\texttt{SATG} Wide-Field Staring/Step-Stare Detection Payload (\texttt{SATG-PLD}):} Uses a three-mirror anastigmat (TMA) telescope ($D_{\text{ap}} \approx 350\text{--}450\text{ mm}$) providing a Wide Field of View (\abbr{WFOV}) covering the full Earth disk ($\sim 17.4^\circ \times 17.4^\circ$) or high-cadence regional step-stare sectors at a frame rate of $\ge 10\text{ Hz}$ and a nadir Ground Sampling Distance (\abbr{GSD}) of $1.0\text{--}2.0\text{ km}$. This achieves a boost-phase launch detection probability of $P_d > 95\%$ (\texttt{R-010}) at a false-alarm rate compatible with $> 99.9\%$ system availability (\texttt{R-015}).
    \item \textbf{\texttt{SATL} Narrow-Field Tracking Payload (\texttt{SATL-PLD}) and Independent Optical Payload \abbr{LOS} Steering Without Affecting Spacecraft Bus Attitude or Other Instrument Pointing Stability (\texttt{R-024}):} Each \texttt{SATL} spacecraft carries a compact Ritchey--Chr\'{e}tien / TMA telescope ($D_{\text{ap}} \approx 200\text{--}250\text{ mm}$, $512 \times 512$ dual-band \abbr{FPA}s) providing a Narrow Field of View (\abbr{NFOV}) of $1.5^\circ \times 1.5^\circ$ (nadir \abbr{GSD} of $25\text{--}50\text{ m}$ at 1,000~km, frame rate $20\text{--}50\text{ Hz}$) steerable across a $\pm 55^\circ$ Field of Regard (\abbr{FOR}). To acquire and lock onto a cued target within \textbf{$< 5.0\text{ seconds}$} while stabilizing the optical Line of Sight (\abbr{LOS}) to \textbf{$< 2\text{ arcsec/s}$ (RMS)} (\texttt{R-024}), \texttt{SATL} steers its payload optical \abbr{LOS} independently via a dual-stage mechanism without slewing the spacecraft bus. Avoiding full-bus slews preserves uninterrupted attitude stability for the platform and maintains continuous pointing lock for the four laser \abbr{OISL} terminals and space-to-ground antennas:
    \begin{enumerate}
        \item \emph{Coarse Stage --- Two-Axis Payload Gimbal / Pointing Mirror Assembly:} A direct-drive brushless DC gimballing mechanism slews the payload optical line of sight at angular rates of $\ge 15^\circ/\text{s}$ and accelerations of $\ge 10^\circ/\text{s}^2$ across the $\pm 55^\circ$ \abbr{FOR}. This brings the target within the $1.5^\circ$ \abbr{NFOV} in $\le 3.0\text{--}3.5\text{ s}$, while spacecraft reaction wheels provide real-time momentum feed-forward compensation so the bus remains stationary ($< 0.05^\circ$ excursion).
        \item \emph{Fine Stage --- Piezo-Actuated Fast Steering Mirror (\abbr{FSM}):} A closed-loop two-axis fine steering mirror at an internal optical pupil image operates with a control bandwidth of $\ge 200\text{ Hz}$ and sub-microradian resolution. Once the coarse gimbal settles, the \abbr{FSM} locks onto the plume centroid within $\le 1.0\text{--}1.5\text{ s}$ (total acquisition $< 5.0\text{ s}$) and actively suppresses micro-vibrations generated by reaction wheels, cryocoolers, and dynamic slewing of the \abbr{SWIR}/\abbr{MWIR} sensor gimbals, directional \abbr{RF} antennas, and agile \abbr{OISL} terminals to hold \abbr{LOS} jitter below $< 2\text{ arcsec/s}$ ($\le 9.7~\mu\mathrm{rad/s}$ RMS). Paired with an optical-bench dual star tracker and fibre-optic gyroscope (\abbr{FOG}) inertial reference unit, each \texttt{SATL} achieves an Attitude/\abbr{LOS} Knowledge Error (\abbr{AKE}) of $\le 10\text{--}15~\mu\mathrm{rad}$ ($\le 2\text{--}3\text{ arcsec}$, $1\sigma$), corresponding to a cross-range target position error of $< 15\text{--}30\text{ m}$ at 1,000--2,000~km slant range.
    \end{enumerate}
    \item \textbf{Stray-Light Rejection, Cryogenic Cooling, and Onboard Proton Despiking (\texttt{R-019}--\texttt{R-023}, \texttt{R-025}):} Both payloads incorporate extended multi-vane blackened optical baffles with knife-edge vanes (\texttt{R-025}) providing $> 10^{-10}$ Point Source Transmittance (PST) attenuation outside a $15^\circ$ solar exclusion angle and $3^\circ$ sunlit Earth-limb exclusion angle. Low-vibration pulse-tube cryocoolers maintain the \abbr{SWIR}/\abbr{MWIR} \abbr{MCT} focal planes at $65\text{ K}\text{--}95\text{ K}$ to suppress dark current. The front-end radiation-tolerant FPGA/SoC signal processing pipeline executes real-time \textbf{proton despiking}---onboard filtering that deletes artificial bright pixel spots (transient ionization spikes) on the infrared detector caused by energetic protons in the space radiation environment---followed by adaptive spatio-temporal background subtraction, track-before-detect filtering (\texttt{R-019}, \texttt{R-022}), and onboard dual-satellite stereoscopic triangulation (\texttt{R-023}).
\end{itemize}

% =============================================================================
% EXTRACTED FROM TRS SECTION 4.4: PRODUCT BREAKDOWN, WBS, SUBSYSTEMS & TABLE 6
% =============================================================================
\subsubsection{4.4 Space Segment Product Breakdown, WBS Alignment, and ``Careful COTS'' Platform Concept}

The physical and functional Product Tree of Configuration Item \texttt{CI-03} (\abbr{SSEG}, \abbr{WBS} \texttt{123456.03}) is aligned with the project Work Breakdown Structure established in Chapter 10 and Section 7.1 of \textbf{AD-01}. Under \textbf{\texttt{123456.03.01} Satellite \abbr{LEO} (\texttt{SATL})} and \textbf{\texttt{123456.03.02} Satellite \abbr{GEO} (\texttt{SATG})}, Chapter 10 of \textbf{AD-01} defines the core WBS branches as \texttt{.01} \emph{Thermal Control} (\texttt{.01.01} Passive, \texttt{.01.02} Active), \texttt{.02} \emph{Attitude Determination and Control System (\abbr{ADCS})}, \texttt{.03} \emph{Power Subsystem (\abbr{EPS})}, \texttt{.04} \emph{Payload (\abbr{PLD})}, \texttt{.05} \emph{Propulsion (\abbr{PROP})}, \texttt{.06} \emph{Avionics} (\texttt{.06.01} Command \& Data Handling, \texttt{.06.02} Communications, \texttt{.06.03} Flight Software), and \texttt{.07} \emph{Constellation Configuration (\texttt{CONF})} (\texttt{.07.01} Phasing Plan, \texttt{.07.02} Orbit Definition). This hierarchy is complemented by two explicit systems-engineering refinements established in Section 7.1 of \textbf{AD-01}:
\begin{enumerate}
    \item \textbf{Explicit Inclusion of the Structures and Mechanisms Subsystem (\abbr{STR/MEC}, \texttt{.08}):} To ensure complete physical accounting alongside WBS elements \texttt{.01} through \texttt{.07}, both \texttt{SATL} and \texttt{SATG} incorporate a dedicated \textbf{Structures and Mechanisms Subsystem (\abbr{STR/MEC}, designated \texttt{123456.03.01.08} and \texttt{123456.03.02.08})} as a primary physical subsystem. \abbr{STR/MEC} encompasses the primary load-bearing CFRP/aluminium-honeycomb bus structure, launch vehicle separation and multi-satellite dispenser rings (supporting 9-satellite-per-plane batch launches for \texttt{SATL}), and solar array and antenna deployment/hold-down release mechanisms. It also includes tuned viscoelastic/fluid micro-vibration isolators that damp mechanical disturbances from reaction wheels, cryocoolers, and dynamic slewing of the \abbr{SWIR}/\abbr{MWIR} sensor gimbals/mirrors, directional \abbr{RF} antennas, and agile \abbr{OISL} terminals (enforcing $< 2\text{ arcsec/s}$ stability per \texttt{R-024}), together with localized Aluminium/Tantalum ($\mathrm{Al}/\mathrm{Ta}$) graded-Z radiation spot-shielding enclosures and external stray-light optical baffles (\texttt{R-025}).
    \item \textbf{Architectural Separation of \abbr{OBDH}, Communications/\abbr{OISL}, and Constellation Configuration (\texttt{CONF}):} Within the \emph{Avionics} (\texttt{.06}) and \emph{Constellation Configuration} (\texttt{.07}) branches of \textbf{AD-01}, this \abbr{TRS} enforces clear functional and hardware subsystem boundaries:
    \begin{itemize}
        \item \textbf{On-Board Data Handling \& Flight Software (\abbr{OBDH}, \texttt{.06.01} \& \texttt{.06.03}):} Encompasses the central Command and Data Handling (\abbr{C\&DH}) flight computers (\texttt{.06.01}), radiation-tolerant FPGA/SoC payload data processors, solid-state mass memory, and onboard Flight Software (\texttt{.06.03}, including \abbr{FDIR}, autonomous mode management, proton despiking, and \texttt{R-023} stereoscopic triangulation).
        \item \textbf{Communications and Mandatory \abbr{OISL} Subsystem (\abbr{TT\&C} \& \abbr{OISL}, \texttt{.06.02}):} Separated as a distinct hardware subsystem encompassing the S/X/Ka-band Space-to-Ground \abbr{RF} transponders, high-gain/omni antennas, cryptographic security modules (\abbr{COMSEC}, \texttt{R-031}), anti-jamming processors (\texttt{R-032}), and the mandatory 1,550~nm laser Optical Inter-Satellite Link (\abbr{OISL}) terminals (\texttt{R-023}, \texttt{R-044}).
        \item \textbf{Constellation Configuration (\texttt{CONF}, \texttt{123456.03.01.07} \& \texttt{123456.03.02.07}):} Treated as a segment-level astrodynamic and mission-design discipline covering the \emph{Phasing Plan} (\texttt{.07.01}) and \emph{Orbit Definition} (\texttt{.07.02}) across the entire \abbr{SSEG} fleet.
    \end{itemize}
\end{enumerate}

To manufacture 72 \texttt{SATL} spacecraft and 6 \texttt{SATG} spacecraft within European programmatic and schedule envelopes while achieving $> 99.9\%$ system availability (\texttt{R-015}) over 7-year (\abbr{LEO}) and 15-year (\abbr{GEO}) lifetimes (\texttt{R-016}), the Space Segment adopts the \textbf{``Careful \abbr{COTS}''} platform philosophy (\textbf{AD-01}, Section 7.2):
\begin{itemize}
    \item \textbf{Selective Radiation Hardening and Spot Shielding:} Rather than procuring expensive, low-throughput QML-V rad-hard components for every board across 72 \abbr{LEO} satellites, \texttt{SATL} uses automotive/industrial-grade high-performance processors and FPGAs up-screened via lot-acceptance radiation testing under the project radiation-assurance rules. These are protected by graded-Z $\mathrm{Al}/\mathrm{Ta}$ spot shielding ($4\text{--}8\text{ mm}$ equivalent $\mathrm{Al}$) and the benign 1,000~km orbit selection. By contrast, the 6 long-life ($15\text{-year}$) \texttt{SATG} geostationary platforms utilize space-grade radiation-hardened avionics for critical bus functions alongside radiation-tolerant high-speed payload processors.
    \item \textbf{Fault-Tolerant Avionics and Redundancy (\texttt{R-017}):} All critical platform chains implement dual cold/warm block redundancy with cross-strapping (\texttt{R-017}), Error Detection and Correction (\abbr{EDAC}) with memory scrubbing, Triple Modular Redundancy (\abbr{TMR}) in FPGA logic, fast-acting Single-Event Latch-up (\abbr{SEL}) current limiters, and autonomous hardware watchdog timers.
\end{itemize}

Table~\ref{tab:sseg_wbs_subsystems} defines the complete Space Segment subsystem breakdown, WBS mapping (matching Chapter 10 of \textbf{AD-01}), functional allocations, and baseline mass and power budgets (explicitly including orbit insertion / transfer propulsion and propellant) for both \texttt{SATL} and \texttt{SATG}.

\begin{table}[H]
\centering
\caption{Space Segment Subsystem Breakdown, WBS Mapping, and Mass/Power Envelopes for \texttt{SATL} and \texttt{SATG}}
\label{tab:sseg_wbs_subsystems}
{\scriptsize\setlength{\tabcolsep}{4pt}\renewcommand{\arraystretch}{1.04}%
\arrayrulecolor{TableBorderBlue}%
\rowcolors{2}{white}{TableRowLightBlue}%
\noindent
\begin{tabularx}{\textwidth}{|>{\raggedright\arraybackslash}p{2.3cm}|>{\raggedright\arraybackslash}p{2.5cm}|>{\raggedright\arraybackslash}X|>{\raggedright\arraybackslash}p{2.3cm}|>{\raggedright\arraybackslash}p{2.3cm}|}
    \hline
    \rowcolor{TableHeaderBlue}
    \thead{Subsystem} & \thead{WBS Code (\texttt{SATL} / \texttt{SATG})} & \thead{Primary Functional \& Hardware Allocation} & \thead{\texttt{SATL} Envelope (1,000~km LEO)} & \thead{\texttt{SATG} Envelope (35,786~km GEO)} \\ \hline
    \textbf{1. Thermal Control (\abbr{TCS})} & \texttt{123456.03.01.01} \newline \texttt{123456.03.02.01} & Passive (\texttt{.01.01}: \abbr{MLI}, radiators) \& Active (\texttt{.01.02}: pulse-tube cryocoolers at 65--95~K, loop heat pipes, heaters). & 16--22~kg \newline 85--125~W & 130--175~kg \newline 550--800~W \\ \hline
    \textbf{2. Attitude Control (\abbr{ADCS})} & \texttt{123456.03.01.02} \newline \texttt{123456.03.02.02} & Dual optical-bench star trackers + \abbr{FOG} IMU ($\text{AKE} \le 10\text{--}15~\mu\mathrm{rad}$); GNSS; 4-wheel skewed reaction wheels; magnetorquers (\texttt{R-024}). & 18--24~kg \newline 45--75~W & 95--130~kg \newline 180--260~W \\ \hline
    \textbf{3. Power Subsystem (\abbr{EPS})} & \texttt{123456.03.01.03} \newline \texttt{123456.03.02.03} & Triple-junction $\mathrm{GaAs}$ deployable solar arrays; $\mathrm{Li\text{-}ion}$ battery pack (35~min \abbr{LEO} / 72~min \abbr{GEO} eclipse); PCDU with \abbr{SEL} limiters. & 32--40~kg \newline 450--650~W (\abbr{EOL}) & 290--370~kg \newline 3.5--5.0~kW (\abbr{EOL}) \\ \hline
    \textbf{4. Mission Payload (\abbr{PLD})} & \texttt{123456.03.01.04} \newline \texttt{123456.03.02.04} & Cooled dual-band \abbr{SWIR}/\abbr{MWIR} $\mathrm{HgCdTe}$ telescope; independent 2-axis gimbal + fine \abbr{FSM} (\texttt{SATL}, $< 5\text{ s}$ lock, $< 2\text{ arcsec/s}$, \texttt{R-009}, \texttt{R-024}). & 48--60~kg \newline 140--190~W & 380--480~kg \newline 950--1,350~W \\ \hline
    \textbf{5. Propulsion (\abbr{PROP})} & \texttt{123456.03.01.05} \newline \texttt{123456.03.02.05} & Orbit insertion/transfer, station-keeping, \abbr{CAM}, \& active \abbr{EOL} disposal (\texttt{SATL}: $\Delta v \ge 380$~m/s; \texttt{SATG}: apogee/\abbr{EOR} transfer to \abbr{GEO} + 15-yr SK + $+300$~km graveyard, \texttt{R-018}). & 14--18~kg (dry) \newline 28--40~kg (prop.) \newline 60--180~W & 140--180~kg (dry) \newline 850--1,150~kg (prop.) \newline 400--1,200~W \\ \hline
    \textbf{6a. Avionics --- \abbr{OBDH} \& \abbr{FSW}} & \texttt{123456.03.01.06} \newline \texttt{123456.03.02.06} \newline (\texttt{.06.01} \& \texttt{.06.03}) & Redundant ``Careful \abbr{COTS}'' / rad-hard \abbr{C\&DH} computers (\texttt{.06.01}, \abbr{TMR}/\abbr{EDAC}) \& Flight Software (\texttt{.06.03}: \abbr{FDIR}, proton despiking \texttt{R-019}, stereo 3D \texttt{R-023}). & 10--14~kg \newline 45--70~W & 45--65~kg \newline 140--210~W \\ \hline
    \textbf{6b. Avionics --- Comms \& \abbr{OISL}} & \texttt{123456.03.01.06} \newline \texttt{123456.03.02.06} \newline (\texttt{.06.02}) & 4~laser \abbr{OISL} terminals/\texttt{SATL} (1,550~nm, $\ge 500$~Mbps, $\ge 6,500$~km); optical \abbr{GEO}--\abbr{LEO} trunks ($\ge 1.5\text{--}2.0$~Gbps, 45,000~km); S/X/Ka \abbr{RF} \abbr{TT\&C} + \abbr{COMSEC} (\texttt{R-031}--\texttt{R-033}, \texttt{R-044}). & 20--25~kg \newline 75--115~W & 130--170~kg \newline 450--700~W \\ \hline
    \textbf{7. Constellation Config.\ (\texttt{CONF})} & \texttt{123456.03.01.07} \newline \texttt{123456.03.02.07} & Phasing Plan (\texttt{.07.01}: Walker Delta $72/8/1$, $\Delta u_{\text{inter}} = 22.5^\circ$ \& $60^\circ$ \abbr{GEO} slots) and Orbit Definition (\texttt{.07.02}: 1,000~km \abbr{LEO} \& 35,786~km \abbr{GEO}, \texttt{R-005}--\texttt{R-008}). & \multicolumn{2}{>{\raggedright\arraybackslash}p{4.6cm}|}{Segment-level astrodynamic design governing all 78 spacecraft.} \\ \hline
    \textbf{8. Structures \& Mech.\ (\abbr{STR/MEC})} & \texttt{123456.03.01.08} \newline \texttt{123456.03.02.08} & Primary CFRP/Al-honeycomb bus; 9-sat dispenser rings (\texttt{SATL}); micro-vibration isolators (wheels, coolers, slewing sensors/antennas/\abbr{OISL}); $\mathrm{Al}/\mathrm{Ta}$ shielding; baffles (\texttt{R-025}). & 42--52~kg \newline 0--15~W & 340--430~kg \newline 0--40~W \\ \hline
    \textbf{Total Spacecraft Envelope} & \textbf{\texttt{SATL} (72x) / \texttt{SATG} (6x)} & \textbf{Complete Flight Unit Mass (incl.\ insertion/transfer propellant \& margin) \& Power Envelope} & \textbf{200--250~kg Dry} \newline \textbf{230--290~kg Wet} \newline \textbf{450--650~W} & \textbf{1,550--1,950~kg Dry} \newline \textbf{2,400--3,100~kg Wet} \newline \textbf{3.5--5.0~kW} \\ \hline
\end{tabularx}}
\end{table}

% =============================================================================
% ARCHIVED FROM TRS ISSUE 1.4 PRIOR TO ISSUE 1.5 REFACTORING
% =============================================================================
\newpage
\section{Archived Space Segment Design Parameters, Orbital Elements, External Interfaces, and Life Profile Envelopes (Extracted from TRS Issue 1.4)}

\subsection{A.1 Baseline Orbital and Astrodynamic Parameters (from TRS Section 4.1)}

To reconcile continuous wide-area boost-phase launch detection with high-precision 3D trajectory estimation (\texttt{R-005}--\texttt{R-008}), the \abbr{EWAS} Space Segment implements a \textbf{dual-tier hybrid orbital architecture} combining a Geostationary Earth Orbit (\abbr{GEO}) detection layer and a high-inclination Low Earth Orbit (\abbr{LEO}) stereoscopic tracking layer:
\begin{enumerate}
    \item \textbf{High-Altitude Detection Layer (\texttt{SATG}, 6~\abbr{GEO} Satellites at 35,786~km, \abbr{WBS} \texttt{123456.03.02}):} Six geostationary spacecraft operate in a circular equatorial orbit ($i = 0^\circ$, semi-major axis $a = 42{,}164.14\text{ km}$, sidereal period $T = 23\text{ h }56\text{ min }4\text{ s}$) spaced uniformly at $60^\circ$ longitude intervals. Viewing an Earth central half-angle of $\lambda_{\text{GEO}} \approx 76.3^\circ$ ($\varepsilon_{\min} = 5^\circ$), the ring provides 24/7 dual-satellite overlapping surveillance across low- and mid-latitude launch corridors up to $\sim 70^\circ$ latitude and continuous cross-layer optical data trunking to European Ground Stations (\abbr{WBS} \texttt{123456.02.05}).
    \item \textbf{Low-Altitude Stereoscopic Tracking Layer (\texttt{SATL}, 72~\abbr{LEO} Satellites at 1,000~km, \abbr{WBS} \texttt{123456.03.01}):} Seventy-two tracking spacecraft operate in a 1,000~km circular orbit ($a = 7{,}378.14\text{ km}$, period $T = 105.12\text{ min}$, $v = 7.35\text{ km/s}$) arranged in a high-inclination ($i = 80^\circ$ nominal within $78^\circ\text{--}82^\circ$) \textbf{Walker Delta $72/8/1$ constellation} comprising $P = 8$ orbital planes ($\Delta\Omega = 45^\circ$ RAAN) with $S = 9$ satellites per plane ($\Delta u = 40^\circ$). Adjacent orbital planes are staggered by a nominal inter-plane phase offset of $\Delta u_{\text{inter}} = 22.5^\circ$ (within the $15.0^\circ\text{--}22.5^\circ$ GMAT-optimized window). This geometry guarantees \textbf{5 to 7 \texttt{SATL} nodes simultaneously in view} ($\varepsilon \ge 5^\circ$) at all times over the Baltic Sea region ($53^\circ\text{N}\text{--}66^\circ\text{N}$) and Arctic launch corridors ($> 60^\circ\text{N}$), ensuring at least one stereoscopic pair with convergence angle $\beta \in [30^\circ, 120^\circ]$.
\end{enumerate}

\begin{table}[H]
\centering
\caption{Baseline Orbital and Astrodynamic Parameters of the EWAS Hybrid Space Segment (Extracted from TRS Issue 1.4 Table 5)}
\label{tab:sseg_orbital_params_archived}
{\footnotesize\setlength{\tabcolsep}{4pt}\renewcommand{\arraystretch}{1.12}%
\arrayrulecolor{TableBorderBlue}%
\rowcolors{2}{white}{TableRowLightBlue}%
\noindent
\begin{tabularx}{\textwidth}{|>{\raggedright\arraybackslash}p{3.6cm}|>{\raggedright\arraybackslash}X|>{\raggedright\arraybackslash}X|}
    \hline
    \rowcolor{TableHeaderBlue}
    \thead{Astrodynamic Parameter} & \thead{SATG --- GEO Detection Layer (123456.03.02)} & \thead{SATL --- LEO Tracking Layer (123456.03.01)} \\ \hline
    \textbf{Primary Operational Role} & Persistent 24/7 boost-phase plume detection (\texttt{R-009}), coarse localization, and cross-layer \abbr{OISL} trunk relay (\texttt{R-042}, \texttt{R-044}). & Agile cued \abbr{NFOV} lock-on ($< 5\text{ s}$, \texttt{R-024}), synchronized stereo tracking, onboard preliminary 3D trajectory (\texttt{R-023}), and \abbr{LOS} package downlink. \\ \hline
    \textbf{Number of Active Satellites} & $N_{\text{GEO}} = 6$ spacecraft (\texttt{R-005}) & $N_{\text{LEO}} = 72$ spacecraft (\texttt{R-005}) \\ \hline
    \textbf{Constellation Topology} & Equatorial ring; uniform $\Delta\lambda = 60^\circ$ longitude slot spacing & Walker Delta $72/8/1$: $P = 8$ planes ($\Delta\Omega = 45^\circ$ RAAN), $S = 9$ sats/plane ($\Delta u = 40^\circ$), $\Delta u_{\text{inter}} = 22.5^\circ$ (\mbox{$15^\circ\text{--}22.5^\circ$} window) \\ \hline
    \textbf{Nominal Altitude ($h$)} & 35,786~km (Geostationary) & 1,000~km (Circular Low Earth Orbit) \\ \hline
    \textbf{Semi-Major Axis ($a$) \& Eccentricity} & $a = 42{,}164.14\text{ km}$; $e \le 0.0005$ & $a = 7{,}378.14\text{ km}$; $e \le 0.001$ \\ \hline
    \textbf{Orbital Inclination ($i$)} & $0.0^\circ \pm 0.05^\circ$ (N--S/E--W station-kept) & $80.0^\circ$ nominal (\mbox{$78^\circ\text{--}82^\circ$} design window) \\ \hline
    \textbf{Sidereal Orbital Period ($T$)} & $23\text{ h }56\text{ min }4.1\text{ s}$ ($1{,}436.07\text{ min}$) & $105.12\text{ min}$ ($6{,}307.2\text{ s}$; $\sim 13.7\text{ orbits/day}$) \\ \hline
    \textbf{Mean Orbital Velocity ($v$)} & $3.075\text{ km/s}$ & $7.350\text{ km/s}$ \\ \hline
    \textbf{Earth Central Half-Angle ($\varepsilon_{\min}=5^\circ$)} & $\lambda_{\text{GEO}} = 76.30^\circ$ (slant range $\le 41{,}120\text{ km}$) & $\lambda_{\text{LEO}} = 25.55^\circ$ (slant range $\le 3{,}290\text{ km}$) \\ \hline
    \textbf{Single-Satellite Earth Coverage} & $38.15\%$ of Earth surface ($\sim 194.6 \times 10^6\text{ km}^2$) & $4.89\%$ of Earth surface ($\sim 24.94 \times 10^6\text{ km}^2$) \\ \hline
    \textbf{Simultaneous Visibility (Baltic $53^\circ\text{--}66^\circ\text{N}$)} & 2--3 \texttt{SATG} nodes continuously visible & 5--7 \texttt{SATL} nodes simultaneously visible 24/7 \\ \hline
    \textbf{Nominal Design Lifetime (\texttt{R-016})} & $\ge 15\text{ years}$ (\abbr{EOL} graveyard $\Delta h \ge +300\text{ km}$) & $\ge 7\text{ years}$ (active controlled/accelerated deorbit) \\ \hline
\end{tabularx}}
\end{table}

\subsection{A.2 Sensing, Tracking, OISL/RF Link, and Flight System Envelopes (from TRS Sections 3.3, 4.2, and 4.3)}

\begin{itemize}
    \item \textbf{Persistent Dual-Band \abbr{SWIR}/\abbr{MWIR} Boost-Phase Detection (\texttt{SATG}):} The 6 \texttt{SATG} nodes perform continuous staring/step-stare surveillance across co-registered \abbr{SWIR} ($2.65\text{--}2.85~\mu\mathrm{m}$) and \abbr{MWIR} ($4.15\text{--}4.45~\mu\mathrm{m}$) bands at $\ge 10.0\text{ Hz}$ (\texttt{R-009}, \texttt{R-020}). Real-time onboard proton despiking and two-colour spectral clutter filtering suppress solar glint, cloud scatter, and non-threat thermal events (\texttt{R-019}, \texttt{R-022}), enabling encrypted \emph{GEO Event Report} downlink within $\le 10.0\text{ s}$ (\texttt{R-010}, \texttt{R-042}).
    \item \textbf{Mandatory $1{,}550\text{ nm}$ Laser \abbr{OISL} Mesh and Cross-Layer Relay Backbone:} Whereas external Space-to-Ground links (\texttt{IF-EXT-01}) use all-weather S/X/Ka-band \abbr{RF}, all Space-to-Space communications use encrypted 1,550~nm laser \abbr{OISL} terminals (\texttt{R-031}--\texttt{R-033}). Each \texttt{SATL} node interconnects via a four-terminal intra-/cross-plane \abbr{LEO} optical mesh ($\ge 500.0\text{ Mbps}$, link range $\ge 6{,}500\text{ km}$, $> 245\text{ km}$ atmospheric clearance) coupled to cross-layer \abbr{GEO}--\abbr{LEO} optical relay trunks ($\ge 1.5\text{--}2.0\text{ Gbps}$, $34{,}800\text{--}45{,}000\text{ km}$). This backbone guarantees sub-3-second delivery of Ground Segment \emph{LEO Tracking Requests} to cued \texttt{SATL} pairs anywhere in orbit (\texttt{R-044}).
    \item \textbf{Cued Stereoscopic Tracking with Independent Optical \abbr{LOS} Steering:} Upon cueing, the two selected \texttt{SATL} spacecraft independently steer their narrow-field payload optical Line of Sight (\abbr{LOS}) across a $\pm 55^\circ$ Field of Regard to lock onto the plume within $< 5.0\text{ s}$ and maintain \abbr{LOS} stability below $< 2.0\text{ arcsec/s}$ (RMS) without slewing the spacecraft bus or disturbing \abbr{OISL} and antenna pointing (\texttt{R-024}). Both \texttt{SATL} nodes track at $\ge 30.0\text{ Hz}$, exchange GNSS-synchronized ($< 50.0~\mu\mathrm{s}$) timestamped \abbr{LOS} unit vectors and $6 \times 6$ covariance matrices over the \abbr{OISL} mesh at $\ge 10.0\text{ Hz}$ to compute an \emph{Onboard Preliminary Trajectory} estimate (\texttt{R-023}), and downlink both \emph{LEO Observation Packages} and the preliminary trajectory to the Ground Segment Data Processing System (\texttt{123456.02.04}) within $\le 5\text{--}7\text{ s}$ (\texttt{R-012}, \texttt{R-026}, \texttt{R-027}, \texttt{R-045}).
    \item \textbf{Flight System Mass and Power Envelopes (from Section 4.3):}
    \begin{itemize}
        \item \textbf{Satellite \abbr{LEO} (\texttt{SATL}, \abbr{WBS} \texttt{123456.03.01}):} Seventy-two (72) 1,000~km tracking satellites (wet launch mass $\le 290.0\text{ kg}$, \abbr{EOL} average power $\ge 550.0\text{ W}$, $\ge 7\text{-year}$ operational lifetime) manufactured under the ``Careful \abbr{COTS}'' electronics philosophy for plane-batch rideshare deployment.
        \item \textbf{Satellite \abbr{GEO} (\texttt{SATG}, \abbr{WBS} \texttt{123456.03.02}):} Six (6) 35,786~km geostationary detection and relay satellites (wet launch mass $\le 3{,}100.0\text{ kg}$, \abbr{EOL} average power $\ge 4{,}200.0\text{ W}$, $\ge 15\text{-year}$ operational lifetime).
    \end{itemize}
\end{itemize}

\begin{table}[H]
\centering
\caption{Summary of Space Segment (\texttt{CI-03}) External Interfaces with Detailed Physical/Band Parameters (Extracted from TRS Issue 1.4 Table 4)}
\label{tab:sseg_external_interfaces_archived}
{\footnotesize\setlength{\tabcolsep}{4pt}\renewcommand{\arraystretch}{1.12}%
\arrayrulecolor{TableBorderBlue}%
\rowcolors{2}{white}{TableRowLightBlue}%
\noindent
\begin{tabularx}{\textwidth}{|>{\raggedright\arraybackslash}p{1.9cm}|>{\raggedright\arraybackslash}p{3.1cm}|>{\raggedright\arraybackslash}p{2.9cm}|>{\raggedright\arraybackslash}X|}
    \hline
    \rowcolor{TableHeaderBlue}
    \thead{Interface ID} & \thead{Interface Name} & \thead{External Counterpart} & \thead{Principal Functional, Physical \& Protocol Characteristics} \\ \hline
    \textbf{IF-EXT-01} & Space-to-Ground \abbr{RF} \abbr{TT\&C} \& Mission Data Interface & Ground Station Network (\abbr{WBS} \texttt{123456.02.05}) \& Ground C2 (\texttt{123456.02.02}) & All-weather bidirectional S-band/X-band \abbr{TT\&C} and high-rate X-band/Ka-band mission data downlinks; \abbr{CCSDS} packet/frame compliant; \abbr{COMSEC} encryption (\texttt{R-031}) and spread-spectrum anti-jamming protection (\texttt{R-032}, \texttt{R-033}). \\ \hline
    \textbf{IF-EXT-02} & Launch Vehicle \& Dispenser Interface & Launch Service Providers (Heavy-Lift \abbr{GEO} \& Medium/Heavy \abbr{LEO}) & Standard mechanical separation rings for \texttt{SATG}; multi-satellite rideshare dispenser mechanisms deploying 9~\texttt{SATL} spacecraft per plane; umbilical power/data/purge lines; launch fairing acoustic, vibrational, and quasi-static envelopes. \\ \hline
    \textbf{IF-EXT-03} & Target Plume \& Background Optical Interface & Boost/Post-Boost Missile Plumes \& Earth/Solar Background & Radiometric entrance aperture interface in \abbr{SWIR} ($\sim 2.7~\mu\mathrm{m}$ $\mathrm{H_2O}/\mathrm{CO_2}$ band) and \abbr{MWIR} ($\sim 4.2\text{--}4.4~\mu\mathrm{m}$ $\mathrm{CO_2}$ red/blue spike bands); solar glint, Earth albedo, cloud-top thermal clutter, and stray-light rejection (\texttt{R-009}, \texttt{R-020}, \texttt{R-025}). \\ \hline
    \textbf{IF-EXT-04} & Space Traffic Management \& Regulatory Interface & Space Surveillance Networks \& \abbr{ITU} Radiocommunication Sector & Reception of Conjunction Data Messages (CDMs) via Ground C2 for collision avoidance at 1,000~km \abbr{LEO} and 35,786~km \abbr{GEO}; compliance with \abbr{ITU} frequency allocations and the project space-debris disposal requirements (\texttt{R-018}). \\ \hline
\end{tabularx}}
\end{table}

\subsection{A.3 Space Segment Life Profile Quantitative Envelopes and Table 6 (from TRS Section 5)}

\begin{itemize}
    \item \textbf{On-Ground AIT, Cleanroom, Storage, and Transport Parameters:} ISO Class~5 dry gaseous nitrogen ($\mathrm{GN_2}$) purged cleanroom conditions ($\mathrm{CVCM} < 0.01\%$, $\mathrm{TML} < 1.0\%$) for cryogenic optical payloads and 1,550~nm \abbr{OISL} assemblies; ISO Class~8 for spacecraft bus integration; up to 24~months spare flight unit storage; dynamic transport loads limited to $\le 3.0\,g$.
    \item \textbf{Launch, LEOP, Bake-Out, Cryocooler, and Commissioning Parameters:} 8 dedicated plane-batch launches of 9 \texttt{SATL} spacecraft each on cylindrical rideshare dispenser stacks ($\le 2{,}800.0\text{ kg}$ per stack) into 1,000~km, $80^\circ$ inclination planes; 6 \texttt{SATG} spacecraft ($\le 3{,}100.0\text{ kg}$ each) on single/dual heavy-lift vehicles; $\ge 3$ independent safety inhibits; post-separation bus tip-off damping from $\le 2.0\text{--}3.0^\circ/\text{s}$ down to $< 0.1^\circ/\text{s}$; 72--96~h outgassing decontamination bake-out at $310\text{--}320\text{ K}$ prior to cooling infrared focal planes to $65\text{--}95\text{ K}$; Walker Delta $72/8/1$ ($\Delta u_{\text{inter}} = 22.5^\circ$) and $60^\circ$ \abbr{GEO} slot insertion.
    \item \textbf{FDIR and End-of-Life Disposal Parameters:} Level~1 sub-100~ms recovery (\abbr{EDAC}, \abbr{TMR}, proton despiking, \abbr{SEL} rail resets); Level~2 redundant unit failover within $\le 2.0\text{ s}$; Level~3 Safe Mode entry within $\le 60.0\text{ s}$; \texttt{SATL} retrograde perigee-lowering deorbit burn $\Delta v_{\text{EOL}} \ge 230.0\text{ m/s}$ ($< 5.0\text{ years}$ demise, $\ge 99.0\%$ reliability, casualty risk $< 10^{-4}$); \texttt{SATG} graveyard re-orbitation $\Delta H \ge +300.0\text{ km}$ ($\Delta v_{\text{EOL}} \ge 15.0\text{ m/s}$).
\end{itemize}

\begin{table}[H]
\centering
\caption{Summary of EWAS Space Segment Life Profile Phases, Operational Modes, and Key Constraints (Extracted from TRS Issue 1.4 Table 6)}
\label{tab:sseg_life_profile_archived}
{\scriptsize\setlength{\tabcolsep}{4pt}\renewcommand{\arraystretch}{1.04}%
\arrayrulecolor{TableBorderBlue}%
\rowcolors{2}{white}{TableRowLightBlue}%
\noindent
\begin{tabularx}{\textwidth}{|>{\raggedright\arraybackslash}p{2.5cm}|>{\raggedright\arraybackslash}p{2.0cm}|>{\raggedright\arraybackslash}p{2.7cm}|>{\raggedright\arraybackslash}X|}
    \hline
    \rowcolor{TableHeaderBlue}
    \thead{Lifecycle Phase / Mode} & \thead{Duration / Timing} & \thead{Applicable CI \& Orbit} & \thead{Space Segment Configuration, Key Activities \& Governing Requirements} \\ \hline
    \textbf{Phase 1: On-Ground \abbr{AIT}, Storage \& Transport} & 18--36~months ($\le 24$~mo storage) & \texttt{EM}, \texttt{STM}, \texttt{QM/PFM}, \texttt{FM} (\texttt{SATL} 72x, \texttt{SATG} 6x) & ISO~5 ($\mathrm{GN_2}$ purge) for \abbr{PLD}/\abbr{OISL} optics; ISO~8 for bus; \abbr{HIL} testing (\texttt{R-039}); \abbr{TVAC}, vibration, shock, and EMC; transport loads $\le 3.0\,g$. \\ \hline
    \textbf{Phase 2: Pre-Launch \& Launch Ascent} & Launch campaign + ascent ($\sim 1\text{--}6\text{ h}$) & 8 plane-launches (9~\texttt{SATL}/launch); \texttt{SATG} \abbr{GTO}/\abbr{GEO} & 9-sat dispenser stack ($\le 2{,}800\text{ kg}$); \texttt{SATG} ($\le 3{,}100\text{ kg}$); \abbr{HDRM}s \& aperture covers locked; $\ge 3$ safety inhibits. \\ \hline
    \textbf{Phase 3: \abbr{LEOP} \& In-Orbit Commissioning} & Weeks 1--6 (\texttt{SATL}) \newline Weeks 1--12 (\texttt{SATG}) & Insertion to 1,000~km ($80^\circ$) \& 35,786~km ($0^\circ$) & Bus detumbling ($\le 2\text{--}3^\circ/\text{s}$ to $< 0.1^\circ/\text{s}$); array deployment; 72--96~h bake-out ($310\text{--}320\text{ K}$); cooler cooldown ($65\text{--}95\text{ K}$); orbit insertion \& Walker Delta phasing ($\Delta u_{\text{inter}} = 22.5^\circ$); \abbr{OISL} mesh validation (\texttt{R-040}). \\ \hline
    \textbf{Phase 4 --- Mode N1: Boost Surveillance} & Continuous 24/7 ($\ge 7\text{ yr}$ \abbr{LEO} / $\ge 15\text{ yr}$ \abbr{GEO}) & \texttt{SATG} primary (6x); \texttt{SATL} high-lat.\ suppl. & Persistent $\ge 10\text{ Hz}$ \abbr{WFOV} \abbr{SWIR}/\abbr{MWIR} staring ($P_d > 95\%$, \texttt{R-009}, \texttt{R-010}); onboard proton despiking \& clutter filtering (\texttt{R-019}, \texttt{R-022}); encrypted \emph{GEO Event Report} downlink in $\le 10.0\text{ s}$ (\texttt{R-042}). \\ \hline
    \textbf{Phase 4 --- Mode N2: Agile Stereo Tracking} & Event-driven ($\le 25\text{--}30\text{ s}$ active window) & 2 cued \texttt{SATL} nodes at 1,000~km + \abbr{OISL} & Triggered by Ground \emph{LEO Tracking Request} ($\le 5\text{ s}$ selection/\abbr{OISL} cueing, \texttt{R-043}, \texttt{R-044}); independent gimbal+\abbr{FSM} lock in $< 5.0\text{ s}$ with $< 2\text{ arcsec/s}$ stability (\texttt{R-024}); \abbr{OISL} \abbr{LOS} vector exchange for \emph{Onboard Preliminary Trajectory} (\texttt{R-023}); dual-product downlink (\texttt{R-011}, \texttt{R-012}, \texttt{R-045}). \\ \hline
    \textbf{Phase 4 --- Mode N3: Mandatory \abbr{OISL} Relay} & Continuous 24/7 background mesh & All 72~\texttt{SATL} \& 6~\texttt{SATG} nodes & 4~laser \abbr{OISL} links/\texttt{SATL} (1,550~nm, $\ge 500$~Mbps, $\ge 6{,}500$~km) + \abbr{GEO}--\abbr{LEO} optical trunks ($\ge 1.5\text{--}2.0$~Gbps, 45,000~km); \abbr{QoS} routing; \abbr{COMSEC} (\texttt{R-031}) \& anti-jamming (\texttt{R-032}, \texttt{R-033}). \\ \hline
    \textbf{Phase 4 --- Mode N4: Station-Keeping \& \abbr{CAM}} & Periodic burns (staggered, 1 node at a time) & 1,000~km \abbr{LEO} ($\pm 0.5^\circ$) \& \abbr{GEO} ($\pm 0.05^\circ$) & Orbit maintenance, momentum dumping, and debris Collision Avoidance Manoeuvres ($P_c > 10^{-4}$, \texttt{IF-EXT-04}); staggered to preserve 24/7 Baltic coverage (\texttt{R-008}, \texttt{R-015}). \\ \hline
    \textbf{Phase 4 --- Mode N5: In-Orbit Calibration} & Scheduled maintenance windows & \texttt{SATL} \& \texttt{SATG} fleet & Radiometric blackbody/deep-space calibration, thermo-elastic boresight updates, and secure uplink of archive-refined algorithms (\texttt{R-004}, \texttt{R-041}). \\ \hline
    \textbf{Phase 5: Contingency (\abbr{FDIR} \& Safe Mode)} & L1 $< 100\text{ ms}$; L2 $< 2\text{ s}$; L3 $\le 60\text{ s}$ entry & Any faulted \texttt{SATL} or \texttt{SATG} node & L1 \abbr{EDAC}/\abbr{TMR}/proton despiking/\abbr{SEL} reset; L2 autonomous A/B redundant unit failover (\texttt{R-017}); L3 Safe Mode closes payload shutter, points arrays to Sun, and activates S-band beacon. \\ \hline
    \textbf{Phase 6: \abbr{EOL} Disposal \& Passivation} & Post-mission ($\le 5\text{ yr}$ decay / immediate \abbr{GEO}) & \texttt{SATL} deorbit; \texttt{SATG} $+300\text{ km}$ graveyard & Active \texttt{SATL} retrograde deorbit ($\Delta v_{\text{EOL}} \ge 230\text{ m/s}$, $\ge 99\%$ reliability); \texttt{SATG} graveyard ($\Delta H \ge +300\text{ km}$); full fluid/battery/RF passivation (\texttt{R-018}). \\ \hline
\end{tabularx}}
\end{table}

\end{document}
